\chapter{Implementation}
\label{chap:implementation}

\section{BoringTun Architecture}

BoringTun~\cite{boringtun} is Cloudflare's userspace WireGuard implementation, written in Rust and deployed on millions of iOS and Android devices along with thousands of Cloudflare Linux servers. It ships as a library crate (\texttt{boringtun}) and a CLI binary (\texttt{boringtun-cli}). The library handles the WireGuard protocol itself but leaves networking and tunnel management to the application layer.

The part of BoringTun that matters for this thesis is the handshake module, which implements the Noise IK pattern. It manages per-peer handshake state machines, performs X25519 operations via the x25519-dalek crate (already a dependency), and uses HMAC-BLAKE2s for chaining key evolution and key derivation. The transport module (ChaCha20-Poly1305 encryption and decryption) is not touched.

Being written in Rust helps here. The ownership and borrowing model rules out buffer overflows, use-after-free, and double-free bugs at compile time, which matters for cryptographic code where such bugs have historically led to exploitable vulnerabilities. It also makes it easy to pull in pure-Rust cryptographic crates from the RustCrypto ecosystem without introducing C FFI boundaries.

\section{Technology Stack}

The implementation relies on the following dependencies:

\textbf{ML-KEM:} The ml-kem crate (version 0.3.0-rc.0) from RustCrypto, a pure-Rust, \texttt{no\_std} implementation of FIPS 203. It supports all three ML-KEM parameter sets and is tested against the NIST Known Answer Test (KAT) vectors. Being pure Rust and \texttt{no\_std} means it runs on constrained targets like the Raspberry Pi without needing any C FFI.

\textbf{X25519:} The x25519-dalek crate is already a dependency of BoringTun and is used for the existing X25519 operations. No changes are required.

\textbf{Key derivation:} HMAC-BLAKE2s is used for chaining key evolution and key derivation, consistent with WireGuard's specification and BoringTun's existing implementation. No separate HKDF or SHA-256 construction is introduced; the hybrid extension reuses the same HMAC-BLAKE2s extraction pattern already present in BoringTun.

\textbf{Test hardware:} Apple MacBook Pro M4 (Apple Silicon, ARM64) and Raspberry Pi 5 (Cortex-A76, ARM64).

\textbf{RNG compatibility:} The ml-kem crate needs \texttt{rand\_core} 0.10, but BoringTun already depends on \texttt{rand\_core} 0.6. Cargo's package renaming feature resolves this: \texttt{rand\_core\_pq} (aliasing \texttt{rand\_core} 0.10) and \texttt{getrandom\_pq} (aliasing \texttt{getrandom} 0.4) are declared as separate optional dependencies, letting both versions coexist in the same binary.

\section{Handshake Module Modifications}

The primary implementation change is to BoringTun's handshake module. All post-quantum code is conditionally compiled behind a Cargo feature flag (\texttt{pq}), ensuring that vanilla builds are completely unaffected: no additional dependencies are compiled, no code paths change, and binary size remains identical. The \texttt{pq} feature gates the ml-kem, rand\_core\_pq, and getrandom\_pq dependencies, as well as all PQ-related message types, state fields, and handshake methods. The following modifications are made when the feature is enabled:

\textbf{State structure extension:} The HandshakeState structure is extended with fields for the ML-KEM ephemeral keypair. On the initiator side, an ml\_kem\_768::\allowbreak EncapsulationKey and ml\_kem\_768::\allowbreak DecapsulationKey are stored. On the responder side, the received encapsulation key (parsed from the initiation message) and the encapsulated ciphertext (to be included in the response message) are stored.

\textbf{Initiation message construction:} When the initiator constructs the handshake initiation message, it generates a fresh ML-KEM-768 keypair (ml\_kem\_768::KeyGen()) and appends the 1,184-byte encapsulation key after the standard WireGuard fields. The encapsulation key is added to the handshake hash h via BLAKE2s mixing, binding it to the Noise transcript.

\textbf{Response message construction:} When the responder processes the initiation message, it extracts the ML-KEM encapsulation key, validates its length, and runs ml\_kem\_768::Encaps() to produce a ciphertext ct and shared secret ss\_kem. The 1,088-byte ciphertext is appended after the standard WireGuard response fields and mixed into the handshake hash.

\textbf{Session key derivation:} At the end of the handshake, both parties mix ss\_kem into the chaining key via one additional HMAC-BLAKE2s extraction step before proceeding to the PSK mixing and final session key derivation. On the initiator side, ss\_kem is obtained by running ml\_kem\_768::Decaps() on the received ciphertext. A new error variant (\texttt{MlKemDecapsulationFailed}) is added to handle decapsulation failures gracefully, causing the handshake to abort without leaking state.

\textbf{New message types:} Two new message types are introduced: type 5 for PQ handshake initiation (1,332 bytes) and type 6 for PQ handshake response (1,180 bytes). The packet parser dispatches on message type and size, routing PQ messages to the corresponding PQ handshake methods. Vanilla builds (without the \texttt{pq} feature) reject these message types with \texttt{InvalidPacket}, ensuring graceful interoperability.

\textbf{MTU handling:} The current implementation relies on IP-level fragmentation for oversized handshake messages. On typical Ethernet paths with a 1,500-byte MTU, both hybrid messages fit in a single UDP datagram. Application-level segmentation and Path MTU Discovery are identified as future work for deployment on paths where IP fragmentation is unreliable.

\section{PSK-Slot Baseline Implementation}

The PSK-slot baseline is implemented as a separate Rust CLI binary, \texttt{pq-psk-tool}, that operates independently of BoringTun. It uses the same ml-kem crate (version 0.3.0-rc.0). The binary provides three subcommands that implement a manual three-step ML-KEM-768 key exchange:

\begin{enumerate}
    \item \texttt{pq-psk-tool keygen}: Generates an ML-KEM-768 keypair. The encapsulation key (1,184 bytes) is written as a base64-encoded file for transfer to the peer. The decapsulation key seed (64 bytes) is written as a separate base64 file and kept secret.
    \item \texttt{pq-psk-tool encaps <ek\_file>}: Reads the peer's encapsulation key, runs ML-KEM.Encaps, and outputs both the ciphertext (as a base64 file for transfer back to the peer) and the 32-byte shared secret (as a hex-encoded PSK file).
    \item \texttt{pq-psk-tool decaps <dk\_file> <ct\_file>}: Reads the local decapsulation key seed and the peer's ciphertext, runs ML-KEM.Decaps, and recovers the identical 32-byte PSK.
\end{enumerate}

The key and ciphertext files are exchanged out-of-band using any secure channel available to the operator (e.g., \texttt{scp}, USB transfer, or a secure messaging channel). The tool itself does not implement a network protocol; it is a purely local utility that reads and writes files. Both peers then configure their WireGuard interface with the derived PSK using \texttt{wg set <iface> peer <PUBKEY> preshared-key psk.hex}.

This is deliberately minimal: no networking code, no daemon, just file I/O. The choice of how to securely move files between peers is left to the operator. For production deployments that need automated PSK distribution, pairing with a daemon like Rosenpass~\cite{rosenpass} or a TLS-protected channel as described by Membrey and Beyel~\cite{membrey2025expressvpn} would make more sense.

\section{Test Harness}

The test infrastructure consists of two components: unit tests within the BoringTun crate and a shell-based integration test script.

\subsection{Unit Tests}

The \texttt{noise::tests} module in BoringTun contains tests for both vanilla and PQ handshake paths, conditionally compiled with \texttt{\#[cfg(\allowbreak feature = "pq")]} and \texttt{\#[cfg(not(\allowbreak feature = "pq"))]}:

\begin{itemize}
\item \textbf{PQ handshake message validation:} Verifies that the PQ initiation message is exactly 1,332 bytes (type 5) and the PQ response is exactly 1,180 bytes (type 6), and that both parse correctly.

\item \textbf{Full PQ handshake:} Exercises the complete four-step exchange (init, response, keepalive, done) and verifies that both peers derive working session keys.

\item \textbf{End-to-end data transport:} Performs a full PQ handshake, then encrypts and decrypts an IPv4 UDP packet, verifying that the decrypted payload matches the original.

\item \textbf{PQ with PSK:} Validates that the hybrid ML-KEM handshake and the traditional PSK mechanism can be used simultaneously, confirming their orthogonality.

\item \textbf{PSK-slot baseline:} Tests the PSK-slot approach by configuring both tunnels with a shared 32-byte PSK (simulating the output of \texttt{pq-psk-tool}), performing a full handshake and data exchange. A companion test verifies that mismatched PSKs cause handshake failure.

\item \textbf{Vanilla rejects PQ:} Verifies that a vanilla build (without the \texttt{pq} feature) rejects PQ message types with \texttt{InvalidPacket}, confirming graceful interoperability.
\end{itemize}

\subsection{Integration Test Script}

A shell script (\texttt{pq-hybrid-test.sh}) automates operational testing on macOS. It provides several modes:

\begin{itemize}
\item \textbf{Demo mode:} Displays PQ vs.\ vanilla packet sizes and runs unit tests for both configurations.

\item \textbf{Local tunnel mode:} Creates two BoringTun peers on loopback interfaces, performs a PQ hybrid handshake, and verifies end-to-end connectivity via ping. Handshake logs are inspected to confirm type 5/6 message exchange.

\item \textbf{Comparison mode:} Runs vanilla and PQ hybrid tunnels side-by-side on separate interfaces and ports, verifying that both configurations achieve successful connectivity.
\end{itemize}

All measurements are taken on the Apple MacBook Pro M4 (local loopback) and on a Raspberry Pi 5 (local loopback and network link between the two machines). Results are compared across all three configurations (vanilla, PSK-slot, hybrid).
